\begin{abstract}
Every new annotation scheme starts with a guideline and almost no data. Supervised parsers need thousands of examples, and an LLM prompted with the guideline ignores its conventions, emits malformed output, and cannot be debugged. We observe that a guideline is mostly a procedure: most steps are deterministic, and only a few call for judgment. We propose \emph{decision orchestration}, which compiles a guideline into an executable procedure. Code runs the steps, enforces well-formedness and assembles the output, and a model is consulted only at decision points, for a schema-constrained fragment or for probabilities over the options of a closed question. We then learn the procedure itself, not model weights, from a handful of gold documents. Decision thresholds are calibrated by measurement, uncertain decisions are routed to a stronger model, and reflection attributes each error to the node that made it and keeps an edit only if it measurably helps. On Chinese Uniform Meaning Representation parsing, orchestration lifts a 9B model from 0 to 17.6 F, two gold documents raise a frontier model from 45.3 to 52.6 F, recalibration on four documents adapts a new decision model (+2.9 F), and routing recovers 75\% of a stronger model's gain at 26\% of its cost. \todo{PARSEME sentence once E9 is run}
\end{abstract}
